% GENERATED FILE. Edit canonical lessons, then run npm run generate:books.
% canonical-source-hash: fnv1a64:920bec5e91e83743
% canonical-lessons: KA-C39-chaha, KA-C39-kaapi, KA-C39-haalu

\chapter{Tea, Coffee, Milk}
\label{ch:ka-tea-coffee-milk}
\hlchaptermodality{39}

\begin{chapteropening}
\textbf{By the end of this chapter:} \emph{I can politely ask for tea, coffee or milk in Kannada, and trace each word along the road it travelled to get here.}

Ask \kn{ದಯವಿಟ್ಟು} \kn{ಚಹಾ}, \kn{ಕಾಫಿ} or \kn{ಹಾಲು}, follow chā through Persian and kāphi through five languages, then explain hālu — the word Chapter 32 already recited unglossed as proof of the p $\to$ h law.
\end{chapteropening}

\section[cahā]{\kn{ಚಹಾ} (cahā) — \textquotedblleft{}tea,\textquotedblright{} which walked to Kannada rather than sailed}

\label{lesson:KA-C39-chaha}



\begin{quote}
Seven body words done. Now the drink you would politely ask for first, in any Karnataka home.
\end{quote}

\subsection*{You'll want to know: \kn{ಚಹಾ}}
\begin{quote}
\textbf{\kn{ಚಹಾ}} — \emph{cahā} — \textbf{tea}
\end{quote}

\begin{quote}
\textbf{\kn{ದಯವಿಟ್ಟು} \kn{ಚಹಾ}.} — \emph{dayaviṭṭu cahā.} — \textquotedblleft{}Tea, please.\textquotedblright{}
\end{quote}

\emph{Dayaviṭṭu} has been yours since Chapter 8.

\begin{cousinweb}
\textbf{\kn{ಚಹಾ}} goes back to Northern Chinese \textbf{chá}, carried overland along Central Asian trade routes into Persian \textbf{\ar{چای}} (\emph{chāy}), and from Persian into Kannada and most of North and South India alike. That is the very same road — Chinese into Persian, Persian into Kannada — that Chapter 20's \textbf{\kn{ಹವಾಮಾನ}} took for its own first half, \emph{hava}.

Tea that instead reached Europe by \textbf{sea} kept a different Chinese form, from the southern coast, which is why English says \textbf{tea} and French says \emph{thé} — one word, two roads, \textquotedblleft{}chai\textquotedblright{} by land and \textquotedblleft{}tea\textquotedblright{} by sea.
\end{cousinweb}

\subsection*{Your turn}
Say these aloud:

\begin{itemize}
\raggedright
  \item \textquotedblleft{}dayaviṭṭu cahā\textquotedblright{} — tea, please
  \item the road — \textquotedblleft{}chá … chāy … cahā\textquotedblright{}
  \item two Persian loans — \textquotedblleft{}hava … cahā\textquotedblright{}
  \item land against sea — \textquotedblleft{}chai … tea\textquotedblright{}
\end{itemize}

\begin{tcolorbox}[breakable,colback=teal!4,colframe=teal!35!black,title={Before you move on}]
Say \textquotedblleft{}tea, please.\textquotedblright{} (\emph{Dayaviṭṭu cahā}.) Name the three languages \emph{cahā} passed through, in order. (\textbf{Chinese}, then \textbf{Persian}, then \textbf{Kannada}.) Which earlier Kannada word took the same Persian road? (\emph{\textbf{\kn{ಹವಾಮಾನ}}}, weather.) Why does English say \textquotedblleft{}tea\textquotedblright{} instead of \textquotedblleft{}chai\textquotedblright{}? (\textbf{It arrived by sea}, carrying a different Chinese form.)
\end{tcolorbox}
\section[kāphi]{\kn{ಕಾಫಿ} (kāphi) — \textquotedblleft{}coffee,\textquotedblright{} the longest chain of loans in this book}

\label{lesson:KA-C39-kaapi}



\begin{quote}
\emph{Cahā} crossed two languages to reach Kannada. This drink crossed four.
\end{quote}

\subsection*{You'll want to know: \kn{ಕಾಫಿ}}
\begin{quote}
\textbf{\kn{ಕಾಫಿ}} — \emph{kāphi} — \textbf{coffee}
\end{quote}

\begin{quote}
\textbf{\kn{ದಯವಿಟ್ಟು} \kn{ಕಾಫಿ}.} — \emph{dayaviṭṭu kāphi.} — \textquotedblleft{}Coffee, please.\textquotedblright{}
\end{quote}

\subsection*{The letters in this word}
\emph{(Skip this if you already read Kannada.)} \textbf{\kn{ಫ}} (\emph{pha}) is new: \textbf{\kn{ಪ}} (\emph{pa}) with an aspirated puff of breath added, the same relationship \textbf{\kn{ಕ}}/\textbf{\kn{ಖ}} and \textbf{\kn{ತ}}/\textbf{\kn{ಥ}} already carry. You met plain \textbf{\kn{ಪ}} back in Chapter 20's \textbf{\kn{ಇಪ್ಪತ್ತು}} (\emph{ippattu}, \textquotedblleft{}twenty\textquotedblright{}) — one of the rare places a native word still opens with \emph{p}, since most word-initial \emph{p}'s became \emph{h}. \textbf{\kn{ಫ}} never went through that law at all: it shows up almost only in borrowed words like this one, holding the foreign sound steady.

\begin{cousinweb}
\textbf{\kn{ಕಾಫಿ}} traces to Arabic \textbf{\ar{قهوة}} (\emph{qahwa}), into Ottoman Turkish \textbf{kahve}, into Dutch \textbf{koffie}, into English \textbf{coffee}, into Kannada \textbf{kāphi} — five languages deep, the longest chain of loans in this book, longer even than \emph{cahā}'s three-language road through China and Persia. South India's filter-coffee tradition is comparatively recent, arriving through colonial-era sea trade — the newer drink, on the older-sounding name.
\end{cousinweb}

\subsection*{Your turn}
Say these aloud:

\begin{itemize}
\raggedright
  \item \textquotedblleft{}dayaviṭṭu kāphi\textquotedblright{} — coffee, please
  \item plain and aspirated — \textquotedblleft{}pa … pha\textquotedblright{}
  \item the chain — \textquotedblleft{}qahwa … kahve … koffie … coffee … kāphi\textquotedblright{}
  \item two drinks, two roads — \textquotedblleft{}cahā, three languages … kāphi, five\textquotedblright{}
\end{itemize}

\begin{tcolorbox}[breakable,colback=teal!4,colframe=teal!35!black,title={Before you move on}]
Say \textquotedblleft{}coffee, please.\textquotedblright{} (\emph{Dayaviṭṭu kāphi}.) Where did you meet plain \textbf{\kn{ಪ}} before, and what makes \textbf{\kn{ಫ}} different? (\textbf{Chapter 20's} \emph{ippattu}; \textbf{\kn{ಫ}} adds a puff of breath and shows up mostly in loanwords.) Name the chain of languages \emph{kāphi} passed through. (\textbf{Arabic, Turkish, Dutch, English, Kannada} — five deep.) Which drink's name travelled further, \emph{cahā}'s or \emph{kāphi}'s? (\emph{\textbf{Kāphi}}'s — five languages against three.)
\end{tcolorbox}
\section[hālu]{\kn{ಹಾಲು} (hālu) — \textquotedblleft{}milk,\textquotedblright{} a word you've already recited without knowing it}

\label{lesson:KA-C39-haalu}



\begin{quote}
Two borrowed drinks in a row. This one is native — and it has been sitting inside a sentence you already said, back in Chapter 32, without anyone telling you what it meant.
\end{quote}

\subsection*{You'll want to know: \kn{ಹಾಲು}}
\begin{quote}
\textbf{\kn{ಹಾಲು}} — \emph{hālu} — \textbf{milk}
\end{quote}

\begin{quote}
\textbf{\kn{ದಯವಿಟ್ಟು} \kn{ಹಾಲು}.} — \emph{dayaviṭṭu hālu.} — \textquotedblleft{}Milk, please.\textquotedblright{}
\end{quote}

\begin{cousinweb}
\textbf{\kn{ಹಾಲು}} continues Old Kannada \emph{pālu}, from Proto-South-Dravidian \textbf{*pāl}, \textquotedblleft{}milk.\textquotedblright{} Tamil keeps the \emph{p} untouched: \textbf{\ta{பால்}} (\emph{pāl}).

Chapter 32 already gave you this exact pair. Teaching \emph{hōgu}, it listed four proofs of the \emph{p} $\to$ \emph{h} law and asked you to say all four aloud:

\begin{quote}
\emph{pōgu … hōgu}, \emph{pattu … hattu}, \emph{\textbf{pāl … hālu}}, \emph{peyar … hesaru}
\end{quote}

You have said \emph{pāl … hālu} before. Nobody told you then that it meant \textquotedblleft{}milk\textquotedblright{} — exactly the gap Chapter 37 found with \emph{bāyi}.

\emph{Hālu} is not only a drink on its own: turmeric stirred into warm milk, \textbf{\kn{ಅರಿಶಿನ} \kn{ಹಾಲು}} (\emph{ariśina hālu}), is an everyday Karnataka drink, its color Chapter 22's other word, \textbf{\kn{ಹಳದಿ}}.
\end{cousinweb}

\subsection*{Your turn}
Say these aloud:

\begin{itemize}
\raggedright
  \item \textquotedblleft{}dayaviṭṭu hālu\textquotedblright{} — milk, please
  \item Chapter 32's four pairs again — \textquotedblleft{}pōgu, hōgu … pattu, hattu … pāl, hālu … peyar, hesaru\textquotedblright{}
  \item turmeric milk, and its color — \textquotedblleft{}ariśina hālu … haḷadi\textquotedblright{}
  \item three drinks now — \textquotedblleft{}cahā, kāphi, hālu\textquotedblright{}
\end{itemize}

\begin{tcolorbox}[breakable,colback=teal!4,colframe=teal!35!black,title={Before you move on}]
Say \textquotedblleft{}milk, please.\textquotedblright{} (\emph{Dayaviṭṭu hālu}.) Where did you already say \emph{pāl … hālu}, before knowing what it meant? (\textbf{Chapter 32's} \emph{hōgu} lesson, reciting the four \emph{p} $\to$ \emph{h} pairs.) Name all four pairs. (\emph{Pōgu}/ \emph{hōgu}, \emph{pattu}/\emph{hattu}, \emph{pāl}/\emph{hālu}, \emph{peyar}/\emph{hesaru}.) What is \emph{ariśina hālu}, and which earlier chapter's color names it? (\textbf{Turmeric milk} — Chapter 22's \textbf{\kn{ಹಳದಿ}}, yellow.)
\end{tcolorbox}
